\section{The need for drainage}
A soil is waterlogged when the water table is so close to the surface that the pores in the root zone are filled with water and the roots cannot breathe. The causes are heavy rain, irrigation with little drainage, seepage from canals and a low permeability of the subsoil. The effects are the lack of oxygen (the roots die within a few days in warm soil), the loss of nitrogen by denitrification, a cold seedbed, and, in dry regions, the rise of salt by capillary action from the water table. Drainage removes the excess water, lowers the water table and allows the leaching of salts.

\section{Kinds of drainage}
\textbf{Surface drainage} removes the water that ponds on the surface, by shaping the land and by open ditches. \textbf{Subsurface drainage} lowers the water table by pipes buried below the surface or by open drains deep enough to reach it. The layouts of subsurface drains are:
\begin{exlist}
\item \kw{Random} (natural) system: pipes in the wet spots only.
\item \kw{Parallel} system: laterals at equal spacing, which feed a collector; for flat, uniform land.
\item \kw{Herringbone}: laterals join the collector at an angle, from both sides; for narrow depressions.
\item \kw{Interceptor} (cut-off) drain: a line across a slope to catch seepage from higher land.
\end{exlist}
The pipes are of corrugated plastic or clay, 60 to 100 mm in diameter for laterals, with slots; they are wrapped in an \kw{envelope} (gravel, geotextile) to prevent fine soil from entering. Spacings of 10 to 30 m and depths of 1.0 to 1.5 m are common.

\section{Steady-state and non-steady-state design}
In the \kw{steady-state} method the recharge (rain or irrigation) is taken as constant, and the water table is assumed to be at a fixed height at midway between drains; the equation of Hooghoudt gives the spacing. In the \kw{non-steady-state} method (Glover and Dumm), the recharge is intermittent and the water table falls after each event, so that the spacing is found from the time in which the water table must fall by a given amount. Non-steady design is more realistic in rain-fed humid areas, steady design in irrigated areas with a constant recharge; the two agree when the recharge is frequent.

\section{Assumptions of the Hooghoudt equation}
The equation is derived with the assumptions of Dupuit and Forchheimer: (1) the flow to the drain is horizontal, and the hydraulic gradient equals the slope of the water table; (2) the soil is homogeneous and isotropic; (3) the drains reach the impermeable layer, or an \emph{equivalent depth} $d$ replaces the actual depth to account for the convergence of the flow near the drain. The equivalent depth is less than the actual depth to the barrier; and when the barrier is deep, the flow below the drain is large, and the radial resistance is the main part of the loss.

\section{Drainage coefficient and the design discharge}
The \kw{drainage coefficient} $q_d$ is the depth of water to be removed per day, in mm\,d$^{-1}$. It is taken from the climate and the crop: 1 to 3 mm\,d$^{-1}$ for irrigated land in arid regions, 5 to 10 for humid regions, and up to 20 for vegetables. A larger coefficient needs closer drains and is more costly; the aim is to keep the water table below the depth that the crop can tolerate for the time it can tolerate it.

\section{Drainage and the environment}
Drainage removes water, but also nitrate, phosphorus and pesticides, and the drain water may harm the water bodies that receive it. Practices to reduce the loss are \emph{controlled drainage} (a weir in the outlet keeps the water table higher between rains), \emph{shallower drains}, the use of buffer strips and wetlands at the outlet, and a better timing of fertiliser. The nitrate concentration of the drain water (Chapter~\ref{chap:17}) is compared with the limit of 11.3 mg N per litre.
